\documentclass[../DoAn.tex]{subfiles}
\usepackage[utf8]{inputenc} 
\usepackage{array} 
\usepackage{graphicx}
\usepackage{caption}
\usepackage{makecell}
\usepackage{multirow}

\begin{document}

% ============================================
% UC01: Học flashcard
% ============================================

\begin{table}[h!]
\centering
\begin{tabular}{|m{5cm}|m{2cm}|m{3cm}|m{3cm}|}
    \hline
    \textbf{Mã Use case} & \textbf{UC01} & \textbf{Tên Use case} & \textbf{Học flashcard} \\
    \hline
    \multicolumn{1}{|m{4cm}|}{\textbf{Tác nhân}} & \multicolumn{3}{|m{9cm}|}{Học sinh (Student)} \\
    \hline
    \multicolumn{1}{|m{4cm}|}{\textbf{Tiền điều kiện}} & \multicolumn{3}{|m{9cm}|}{Học sinh đã đăng nhập vào hệ thống và có quyền truy cập vào bộ flashcard (là chủ sở hữu hoặc thành viên lớp học có bộ flashcard đó)} \\
    \hline
    \multirow{10}{*}{\textbf{Luồng sự kiện chính}} & 
    \multicolumn{3}{|m{9cm}|}{
    \begin{tabular}{|m{1cm}|m{2cm}|m{4.7cm}|} 
        \hline
        \textbf{STT} & \textbf{Thực hiện bởi} & \textbf{Hành động} \\
        \hline
        1 & Học sinh & Chọn bộ flashcard cần học từ danh sách \\
        \hline
        2 & Hệ thống & Tạo phiên học mới và lấy flashcard tiếp theo theo thứ tự ưu tiên (due terms → new terms → reviewed terms) \\
        \hline
        3 & Hệ thống & Hiển thị mặt trước của flashcard \\
        \hline
        4 & Học sinh & Xem mặt trước, suy nghĩ và lật xem mặt sau \\
        \hline
        5 & Học sinh & Đánh giá mức độ nhớ (recall score từ 0-5) \\
        \hline
        6 & Hệ thống & Tính toán EF, interval, next\_review\_date bằng thuật toán SM-2 và lưu StudyActivity \\
        \hline
        7 & Hệ thống & Cập nhật ProgressSummary (mastered\_terms, reviewing\_terms, forgotten\_terms, completion\_rate) \\
        \hline
        8 & Học sinh & Lặp lại bước 2-7 cho đến khi hết flashcard hoặc kết thúc phiên học \\
        \hline
    \end{tabular}} \\
    \hline
    \multirow{4}{*}{\textbf{Luồng sự kiện thay thế}} & 
    \multicolumn{3}{|m{9cm}|}{
    \begin{tabular}{|m{1cm}|m{2cm}|m{4.7cm}|} 
        \hline
        \textbf{STT} & \textbf{Thực hiện bởi} & \textbf{Hành động} \\
        \hline
        3a & Hệ thống & Thông báo lỗi nếu bộ flashcard không tồn tại hoặc không có quyền truy cập \\
        \hline
        3b & Hệ thống & Thông báo nếu bộ flashcard trống (không có flashcard nào) \\
        \hline
        8a & Hệ thống & Thông báo lỗi nếu không thể lưu kết quả đánh giá vào cơ sở dữ liệu \\
        \hline
        5a & Học sinh & Yêu cầu AI tạo đoạn văn minh họa hoặc test từ flashcard hiện tại \\
        \hline
    \end{tabular}} \\
    \hline
    \multicolumn{1}{|m{4cm}|}{\textbf{Hậu điều kiện}} & \multicolumn{3}{|m{9cm}|}{Phiên học được lưu thành công, tiến độ học tập được cập nhật, thời gian ôn lại tiếp theo được tính toán cho từng flashcard} \\
    \hline
\end{tabular}
\caption{Thông tin chi tiết về Use Case Học flashcard}
\end{table}

\vspace{0.5cm}

\begin{table}[h!]
\centering
\begin{tabular}{|m{1cm}|m{2.5cm}|m{3cm}|m{1cm}|m{2.5cm}|m{3.5cm}|}
    \hline
    \multicolumn{6}{|c|}{\textbf{Dữ liệu đầu vào}} \\
    \hline
    \textbf{STT} & \textbf{Trường dữ liệu} & \textbf{Mô tả} & \textbf{Bắt buộc} & \textbf{Điều kiện hợp lệ} & \textbf{Ví dụ} \\
    \hline
    1 & studyset\_id & Mã định danh của bộ flashcard & Có & UUID hợp lệ & "123e4567-e89b-12d3-a456-426614174000" \\
    \hline
    2 & session\_type & Loại phiên học & Có & "flashcard", "quick\_review" & "flashcard" \\
    \hline
    3 & recall\_score & Điểm đánh giá mức độ nhớ & Có & Số nguyên từ 0 đến 5 & 4 \\
    \hline
    4 & response\_time & Thời gian phản hồi (giây) & Không & Số thực dương & 3.5 \\
    \hline
    5 & hint\_used & Có sử dụng gợi ý hay không & Không & Boolean & true \\
    \hline
    6 & mode & Chế độ học & Không & "random" hoặc null & "random" \\
    \hline
\end{tabular}
\caption{Dữ liệu đầu vào cho Use Case Học flashcard}
\end{table}

\vspace{0.5cm}

\begin{table}[h!]
\centering
\begin{tabular}{|m{1cm}|m{2.5cm}|m{3cm}|m{1cm}|m{2.5cm}|m{3.5cm}|}
    \hline
    \multicolumn{6}{|c|}{\textbf{Dữ liệu đầu ra}} \\
    \hline
    \textbf{STT} & \textbf{Trường dữ liệu} & \textbf{Mô tả} & \textbf{Bắt buộc} & \textbf{Điều kiện hợp lệ} & \textbf{Ví dụ} \\
    \hline
    1 & session\_id & Mã định danh phiên học & Có & UUID hợp lệ & "123e4567-e89b-12d3-a456-426614174000" \\
    \hline
    2 & terms & Danh sách flashcard trong phiên học & Có & Mảng các đối tượng term & [\{term\_id, term\_text, definition\}] \\
    \hline
    3 & next\_review\_date & Ngày ôn lại tiếp theo & Có & Định dạng ngày hợp lệ & "2025-01-15T10:30:00" \\
    \hline
    4 & ef & Hệ số dễ dàng (Ease Factor) & Có & Số thực từ 1.3 đến 2.5 & 2.1 \\
    \hline
    5 & interval & Khoảng thời gian ôn lại (ngày) & Có & Số nguyên dương & 7 \\
    \hline
    6 & session\_summary & Tóm tắt kết quả phiên học & Có & Đối tượng chứa thống kê & \{total\_reviewed: 20, average\_score: 3.5\} \\
    \hline
    7 & Thông báo lỗi (nếu có) & Lý do thất bại & Không & Chỉ hiển thị khi có lỗi & "Không tìm thấy bộ flashcard" \\
    \hline
\end{tabular}
\caption{Dữ liệu đầu ra cho Use Case Học flashcard}
\end{table}

\end{document}
